Práca sa zaoberá automatickou tvorbou znalostných grafov zo slovenských právnych predpisov pomocou veľkých jazykových modelov. Klasický vektorový RAG nad zákonmi zlyháva z troch dôvodov: paragrafy sa cez celý dokument odkazujú navzájom, tabuľky a vzorce sa v plochej textovej konverzii rozpadnú.

V práci je navrhnutá a implementovaná pipeline, ktorá rieši všetky tri problémy. Predspracovanie oddelene extrahuje tabuľky a vzorce. Trojkroková extrakcia najprv v otvorenej doméne zhromaždí typy entít a vzťahov v dokumente, potom ich pomocou väčšieho modelu zjednotí do koherentnej schémy, a nakoniec touto schémou riadi konkrétnu extrakciu inštancií. Výsledné uzly a hrany sa vkladajú do grafovej databázy Neo4j s troma vektorovými indexmi nad chunkami, entitami a vzťahmi pre hybridné vyhľadávanie.

Riešenie je vyhodnotené nad zákonom č. 222/2004 Z. z. o dani z pridanej hodnoty. Štyri kombinácie modelu a režimu promptovania sú porovnané voči ručne anotovanej referencii pomocou F1 skóre, miery halucinácií, opomenutí a normalizovanej grafovej editovacej vzdialenosti. Druhá časť experimentov porovnáva odpovede klasického vektorového RAG-u s GraphRAG-om nad rovnakým zákonom a hodnotí, či navrhnutý prístup presnejšie nachádza podporujúce ustanovenia.
